% !TEX program = xelatex
\documentclass[12pt,a4paper]{article}
\usepackage{fontspec}
\usepackage{polyglossia}
\setmainlanguage{russian}
\setotherlanguage{english}
\setmainfont{DejaVu Serif}
\setsansfont{DejaVu Sans}
\setmonofont{DejaVu Sans Mono}
\usepackage{geometry}
\geometry{margin=2.5cm}
\usepackage{amsmath,amssymb,amsthm,mathtools,bm}
\usepackage{microtype}
\usepackage{enumitem}
\usepackage{hyperref}
\hypersetup{colorlinks=true, linkcolor=blue, citecolor=blue, urlcolor=blue}
\newtheorem{definition}{Определение}
\newtheorem{axiom}{Аксиома}
\newtheorem{theorem}{Теорема}
\newtheorem{proposition}{Утверждение}
\newtheorem{remark}{Замечание}
\DeclareMathOperator{\Valid}{ValidInterface}
\DeclareMathOperator{\Interface}{Interface}
\DeclareMathOperator{\Auto}{Auto}
\DeclareMathOperator{\AlwaysAct}{AlwaysAct}
\DeclareMathOperator{\Input}{Input}
\DeclareMathOperator{\Adm}{Adm}
\DeclareMathOperator{\Dir}{Dir}
\DeclareMathOperator{\AutoAct}{AutoAct}
\DeclareMathOperator{\PassiveChannel}{PassiveChannel}
\DeclareMathOperator{\Control}{Control}
\DeclareMathOperator{\Der}{Der}
\DeclareMathOperator{\Struct}{Struct}
\DeclareMathOperator{\Dist}{Dist}
\DeclareMathOperator{\End}{End}
\DeclareMathOperator{\Tr}{Tr}
\DeclareMathOperator{\Attr}{Attr}
\DeclareMathOperator{\Coh}{Coh}
\DeclareMathOperator{\Sync}{Sync}
\DeclareMathOperator{\Loss}{Loss}
\DeclareMathOperator{\Stab}{Stab}
\newcommand{\M}{\mathcal M}
\newcommand{\B}{\mathcal B}
\newcommand{\Hcal}{\mathcal H}
\newcommand{\Ispace}{\mathfrak I}
\newcommand{\XiAuto}{\Xi_I^{\mathrm{auto}}}
\newcommand{\XiMan}{\mathcal M_{\Xi}^{\mathrm{auto}}}
\newcommand{\NXi}{\mathcal N_{\Xi}}
\newcommand{\PiAmp}{\Pi_{\mathrm{amp}}}
\newcommand{\PiStr}{\Pi_{\mathrm{str}}}
\newcommand{\PiStruct}{\Pi_{\mathrm{struct}}}

\title{\textbf{Автоматически структурированные интерфейсные поля в TMI-Physics 2.2}\\
\large От квантовых распределений возможностей к масштабируемым интерфейсным многообразиям}
\author{Salkutsan Aleksey}
\date{2026}
\begin{document}
\maketitle
\begin{abstract}
Работа формулирует версию TMI-Physics 2.2 как расширение минимального физического слоя Теории многообразных интерфейсов. В TMI-Physics 2.1 минимальное интерфейсное поле задавалось скаляром \(\chi_I\). В версии 2.2 это поле сохраняется как амплитудная проекция полного автоматически структурированного интерфейсного поля \(\Xi_I^{\mathrm{auto}}\). Новая версия связывает геометрическую допустимость, квантовое распределение возможностей, аксиому первичного интерфейса, дефиниционную автоматичность интерфейса, автоструктурирование поля и масштабирование к многообразию интерфейсных полей.
\end{abstract}
\tableofcontents
\newpage

\section{Статус версии}
TMI-Physics 2.2 не отменяет TMI-Physics 2.1. Минимальное поле \(\chi_I\) остаётся рабочей скалярной проекцией:
\begin{equation}
\boxed{\chi_I=\PiAmp(\XiAuto).}
\end{equation}
Таким образом, новая версия является расширением, а не заменой минимальной модели.

\section{Минимальное ядро TMI-Physics 2.1}
Минимальная цепочка имела вид:
\begin{equation}
\boxed{g\Rightarrow\Adm_c(g)\Rightarrow\Hcal_I(g)\Rightarrow\Psi_I\Rightarrow\Dist_{\mathrm{poss}}(\Psi_I;g)\Rightarrow\chi_I\Rightarrow\Struct_I.}
\end{equation}
Смысл: геометрия задаёт допустимость; квантовая теория задаёт распределение; интерфейсное поле задаёт структурирование.

\section{Расширенное ядро TMI-Physics 2.2}
Версия 2.2 вводит полную форму:
\begin{equation}
\boxed{g\Rightarrow\Adm_c(g)\Rightarrow\Hcal_I(g)\Rightarrow\Psi_I\Rightarrow\Dist_{\mathrm{poss}}(\Psi_I;g)\land I_0\Rightarrow\XiAuto\Rightarrow\XiMan\Rightarrow\Struct_I(\M).}
\end{equation}
Здесь \(I_0\) не является событием после геометрии во времени. Он является онтологическим основанием различия, без которого невозможны сторона, переход, отношение и интерпретация.

\section{Геометрическая допустимость}
Пусть \((\M,g)\) - пространство-время с метрикой \(g\). Область причинно допустимых переходов:
\begin{equation}
\Adm_c(g):=\{(x,y)\in\M\times\M\mid y\in J_g^+(x)\}.
\end{equation}
Первый слой:
\begin{equation}
\boxed{g\Rightarrow\Adm_c(g).}
\end{equation}

\section{Квантовое распределение возможностей}
На области допустимых переходов строится гильбертово пространство:
\begin{equation}
\Hcal_I(g):=L^2(\Adm_c(g),\B_g,\mu_g).
\end{equation}
Если \(\Psi_I\in\Hcal_I(g)\), то состояние задаёт распределение возможностей:
\begin{equation}
\boxed{\Psi_I\Rightarrow\Dist_{\mathrm{poss}}(\Psi_I;g).}
\end{equation}

\section{Аксиоматический источник автоматичности}
Из TMI-Core 0.2:
\begin{equation}
\boxed{N_0\xrightarrow{I_0}\Delta_0\to(A,\neg A).}
\end{equation}
\begin{equation}
\boxed{\Valid(I)\Rightarrow\Auto(I).}
\end{equation}
\begin{equation}
\boxed{\Auto(I)\equiv(\Input(I)\land\Adm(I)\land\Dir(I)\Rightarrow\AutoAct(I)).}
\end{equation}
Если физическое поле является интерфейсным полем, то:
\begin{equation}
\boxed{\Xi_I\in\Ispace_{\mathrm{phys}}\Rightarrow\Auto(\Xi_I)\Rightarrow\XiAuto.}
\end{equation}

\section{Автоматически структурированное интерфейсное поле}
\begin{definition}
Автоматически структурированное интерфейсное поле - операторное поле
\begin{equation}
\XiAuto\in\Gamma(\End(\Hcal_I)),
\end{equation}
обладающее внутренним законом самообновления:
\begin{equation}
\partial_\tau\XiAuto=\mathcal A_I(\XiAuto,g,\Psi_I,\Adm_c(g),\Dist_{\mathrm{poss}}(\Psi_I;g)).
\end{equation}
\end{definition}

\section{Внутренняя теорема структурирования}
\begin{theorem}
Если \(\XiAuto\in\Gamma(\End(\Hcal_I))\), его динамика автономна, замкнута в допустимом пространстве и существует функционал устойчивости \(\mathcal L_I\geq0\), такой что
\begin{equation}
\frac{d}{d\tau}\mathcal L_I[\XiAuto]\leq0,
\end{equation}
то
\begin{equation}
\boxed{\XiAuto\Rightarrow\Struct_I.}
\end{equation}
\end{theorem}

\begin{proof}
Поле \(\XiAuto\) действует на пространство возможностей \(\Hcal_I\). Автономность исключает внешний управляющий оператор, замкнутость сохраняет допустимость поля, а невозрастание \(\mathcal L_I\) задаёт движение к устойчивому режиму. Этот режим обозначается \(\Struct_I:=\Attr(\mathcal A_I)\). Следовательно, \(\XiAuto\Rightarrow\Struct_I\).
\end{proof}

\section{Проекции}
\begin{equation}
\boxed{\XiAuto\to\Xi_I\to\chi_I.}
\end{equation}
\begin{equation}
\boxed{\chi_I^2(x)=\frac{1}{N_I}\Tr_I\left[(\XiAuto(x))^\dagger\XiAuto(x)\right].}
\end{equation}
\begin{equation}
\boxed{\Struct_I=\mathcal S_I(g,\Psi_I,\XiAuto,\Adm_c(g),\Dist_{\mathrm{poss}}(\Psi_I;g)).}
\end{equation}

\section{Масштабирование}
\begin{equation}
\boxed{\XiAuto\to\{\Xi_{I,k}^{\mathrm{auto}}\}_{k\in K}\to\NXi\to\XiMan\to\Struct_I(\M).}
\end{equation}
Многообразие автоматически структурированных интерфейсных полей:
\begin{equation}
\boxed{\XiMan=(V_\Xi,E_\Xi,O_\Xi,\mathcal A_\Xi,\PiStr,\PiAmp,\PiStruct).}
\end{equation}

\section{Согласованность и синхронизация}
Для связанных полей:
\begin{equation}
\Coh_\Xi(k,l)=1-\Loss_\Xi(\Xi_{I,k}^{\mathrm{auto}},\Xi_{I,l}^{\mathrm{auto}}).
\end{equation}
Условие многообразия:
\begin{equation}
\Coh_\Xi(\XiMan)\to1,\qquad \Sync_\Xi(\XiMan)\to1.
\end{equation}

\section{Статус утверждений}
\begin{itemize}
\item Определения: \(\Valid(I)\Rightarrow\Auto(I)\).
\item Внутренние следствия: \(\XiAuto\land\Stab(\XiAuto)\Rightarrow\Struct_I\).
\item Физические гипотезы: \(\XiAuto\) существует как физический механизм.
\item Эксперименты: \(T_2^{TMI}>T_2^{base}\) в квантовой реализации.
\end{itemize}

\section{Заключение}
TMI-Physics 2.2 формулирует физический слой как проекцию обновлённого ядра автоматических интерфейсов. Скалярное поле \(\chi_I\) сохраняется как минимальная наблюдаемая проекция, а полная форма \(\XiAuto\) задаёт автоматическое структурирование распределённых возможностей и масштабирование к интерфейсному многообразию.

\begin{thebibliography}{9}
\bibitem{TMIcore} Salkutsan Aleksey. \textit{Theory of Manifold Interfaces}. \href{https://doi.org/10.5281/zenodo.19940675}{doi:10.5281/zenodo.19940675}.
\bibitem{TMIphysics21} Salkutsan Aleksey. \textit{Theory of Manifold Interfaces: A Regularized Effective Physical Layer, Observable Consequences, and Limitations}. 2026.
\end{thebibliography}
\end{document}
